\chapter{\ifproject%
\ifenglish Experimentation and Results\else การทดลองและผลลัพธ์\fi
\else%
\ifenglish System Evaluation\else การประเมินระบบ\fi
\fi}

\section{การตั้งค่าการทดลอง (Experimental Setup)}

การทดลองดำเนินการในพื้นที่แปลงฟื้นฟูป่าโปงแยงนอก จังหวัดเชียงใหม่ 
ซึ่งมีลักษณะเป็นพื้นที่ป่าไม้ที่มีความหนาแน่นของพืชพรรณในระดับปานกลางถึงสูง 
เหมาะสำหรับการทดสอบระบบวัดระยะในสภาพแวดล้อมจริง

ในการทดลองนี้เลือกต้นไม้เป้าหมายจำนวน 3 ต้น 
เพื่อใช้เป็นตัวอย่างในการประเมินความแม่นยำของระบบ

โดรน DJI M30T ถูกตั้งค่าให้ลอยตัวนิ่ง (Hovering) เหนือยอดไม้ 
เพื่อให้สามารถวัดระยะในแนวดิ่งจากเซ็นเซอร์ VPS ได้อย่างต่อเนื่อง 
โดยในระหว่างการทดลอง โดรนจะรักษาระดับความสูงคงที่ 
และเคลื่อนที่ในแนวราบเพื่อทำการวัดระยะด้วยระบบ UWB

อุปกรณ์ที่ใช้ในการทดลองประกอบด้วย:
\begin{itemize}
    \item โมดูล ESP32 UWB (Tag และ Anchor)
    \item โดรน DJI M30T
    \item คอมพิวเตอร์สำหรับบันทึกข้อมูล (CSV Logging)
    \item ซอฟต์แวร์ WebODM สำหรับสร้างแบบจำลองสามมิติ
\end{itemize}

\section{วิธีการเก็บข้อมูล (Data Collection)}

กระบวนการเก็บข้อมูลแบ่งออกเป็น 2 ส่วนหลัก ได้แก่ 
(1) การเก็บข้อมูลระยะทางจากระบบ UWB และเซ็นเซอร์ของโดรน 
และ (2) การเก็บภาพถ่ายเพื่อสร้างแบบจำลองสามมิติ

\subsection{การเก็บข้อมูลระยะทาง (UWB และ VPS)}

หลังจากทำการติดตั้งอุปกรณ์ UWB โดยติดตั้ง Tag บนโดรน 
และ Anchor บนพื้นดินบริเวณโคนต้นไม้เป้าหมาย 
ดำเนินการวัดค่าระยะทางตามขั้นตอนดังนี้

\begin{enumerate}
    \item โดรนบินขึ้นไปลอยตัวเหนือยอดไม้ 
    เพื่อให้เซ็นเซอร์ VPS วัดระยะจากโดรนถึงยอดไม้

    \item โดรนเคลื่อนที่ในแนวราบโดยรักษาระดับความสูงเดิม 
    เพื่อวัดระยะระหว่าง Tag และ Anchor ด้วยระบบ UWB

    \item ในบางช่วงพบการสูญเสียสัญญาณ 
    เมื่อโดรนอยู่เหนือยอดไม้โดยตรง 
    ซึ่งเกิดจากสภาวะ Non-Line-of-Sight (NLOS)

    \item ข้อมูลระยะทางจาก UWB และ VPS 
    ถูกบันทึกในรูปแบบไฟล์ CSV เพื่อนำไปวิเคราะห์
\end{enumerate}

\subsection{การเก็บข้อมูลภาพถ่าย (Image Acquisition)}

ก่อนการวัดระยะด้วย UWB 
ได้ทำการใช้โดรนบินสำรวจพื้นที่และเก็บภาพถ่ายทางอากาศ 
โดยมีการซ้อนทับของภาพ (Overlap) เพียงพอสำหรับการประมวลผล

ภาพถ่ายที่ได้ถูกนำไปประมวลผลด้วยซอฟต์แวร์ WebODM 
เพื่อสร้างแบบจำลองสามมิติในรูปแบบ 3D Point Cloud 
ซึ่งใช้เป็นค่ามาตรฐานอ้างอิงในการวัดความสูงของต้นไม้

\section{วิธีการคำนวณ (Computation)}

การคำนวณความสูงของต้นไม้ในงานทดลองนี้ 
อาศัยสมการที่ \eqref{eq:tree_height} 
โดยใช้ค่าระยะทางจากระบบ UWB ($D_{UWB}$) 
และค่าระยะจากเซ็นเซอร์ VPS ($D_{VPS}$) 
เพื่อนำมาคำนวณความสูงของต้นไม้

สำหรับค่าความสูงอ้างอิง 
ได้มาจากการวัดโดยตรงจากแบบจำลองสามมิติ (3D Point Cloud) 
ที่สร้างขึ้นจากซอฟต์แวร์ WebODM 
โดยการเลือกจุดยอดไม้และจุดระดับพื้นดิน 
และคำนวณระยะในแนวดิ่งระหว่างสองจุดดังกล่าว

\section{ผลการทดลอง (Results)}


\section{การเปรียบเทียบผล (Comparison)}


\section{การวิเคราะห์ผล (Discussion)}

